% Ad Familiares 11.25 --- headnote
% ad-familiares-11-25, 18 June 43 BC, Rome

\textit{Cicero to D. Junius Brutus Albinus, from Rome on 18
June 43 BC --- Perseus dateline \textit{Scr. Romae xiv K.
Quint. a. 711 (43)}. Lupus, a courier already shuttling
between Rome and Brutus's army, has appeared at Cicero's
door asking for a dispatch on the spot. The senate's
official gazette is going out separately; Cicero, with
nothing fresh to report, makes the letter a study in
restraint.}

\textit{The note is light in tone but freighted with anxiety.
``In you and your colleague'' --- D. Brutus and Plancus,
consuls-designate for 42 --- rests all hope. About M. Junius
Brutus in Macedonia there is still no firm news, though
Cicero is following D. Brutus's standing instruction and
writing to him privately to bring his army into the common
cause. The ``inward evil of the city'' is no small one ---
the unnamed reference is to Octavian's increasingly open
manoeuvring in Rome through that month, which would come
to a head in his march on the city in late July. Cicero
catches himself spilling onto a second page in spite of his
correspondent's much-vaunted Laconic brevity, signs off
\textit{vince et vale}, and dates the letter to the day
fourteen before the Kalends of Quintilis (18 June).}
